\documentclass[11pt,a4paper]{article}
\usepackage[italian]{babel}
\usepackage[margin=2.5cm]{geometry}
\usepackage{parskip}
\usepackage{amsmath, amssymb, amsthm}
\usepackage{booktabs}
\usepackage{array}
\usepackage{xcolor}
\usepackage{needspace}
\usepackage{enumitem}
\usepackage{listings}
\usepackage{tikz}
\usepackage[most]{tcolorbox}
\usepackage[hidelinks]{hyperref}
\usetikzlibrary{decorations.pathreplacing, shapes.geometric, arrows.meta, positioning}

% Niente vedove/orfane: i paragrafi non lasciano righe isolate a fine/inizio pagina
\widowpenalty=10000
\clubpenalty=10000

% Elenchi più compatti
\setlist{itemsep=2pt, parsep=0pt, topsep=4pt, partopsep=0pt}

% --- Riquadri con bordo nero, senza numero (si spezzano tra due pagine) ---
% Uso: \begin{definizione} ... \end{definizione}
%      \begin{teorema}[Nome] ... \end{teorema}  (il nome è facoltativo)
\tcbset{nota/.style={enhanced jigsaw, breakable, colback=white,
  colframe=black, boxrule=0.5pt, sharp corners}}
\NewTColorBox{definizione}{}{nota,
  before upper={\textbf{Definizione.}\ }}
\NewTColorBox{teorema}{o}{nota,
  before upper={\textbf{Teorema\IfValueT{#1}{ (#1)}.}\ }}
\NewTColorBox{proposizione}{o}{nota,
  before upper={\textbf{Proposizione\IfValueT{#1}{ (#1)}.}\ }}
\NewTColorBox{proprieta}{o}{nota,
  before upper={\textbf{Proprietà\IfValueT{#1}{ (#1)}.}\ }}
\NewTColorBox{assioma}{o}{nota, boxrule=1pt,
  before upper={\textbf{Assioma\IfValueT{#1}{ (#1)}.}\ }}

% --- Ambienti semplici (senza riquadro) ---
% Il titolo sta in cima al blocco, anche se il contenuto inizia con un riquadro o
% con codice; e non resta mai da solo in fondo a una pagina.
% Uso: \begin{esempio}[Titolo facoltativo] ... \end{esempio}
\newcommand{\newrunin}[2]{%
  \NewDocumentEnvironment{#1}{o}{%
    \par\Needspace{4\baselineskip}\addvspace{\medskipamount}\noindent
    \textbf{#2}\IfValueT{##1}{ (##1)}\textbf{.}\ \ignorespaces}%
  {\par\addvspace{\medskipamount}}}
\newrunin{esempio}{Esempio}
\newrunin{esercizio}{Esercizio}
\newrunin{osservazione}{Osservazione}

% V e F nelle tavole di verità
\newcommand{\V}{\textsf{V}}
\newcommand{\F}{\textsf{F}}

% --- Codice C ---
% Uso: \begin{codice} ... \end{codice}      codice C numerato
%      \begin{comando} ... \end{comando}    comandi da terminale
%      \begin{risultato} ... \end{risultato} output di un programma
%      \lstinline|printf("ciao");|           codice dentro al testo
\lstdefinestyle{c}{
  language=C,
  basicstyle=\ttfamily\small,
  keywordstyle=\color{blue}\bfseries,
  commentstyle=\color{gray}\itshape,
  stringstyle=\color{red!70!black},
  numbers=left,
  numberstyle=\tiny\color{gray},
  numbersep=8pt,
  columns=flexible,
  keepspaces=true,
  breaklines=true,
  showstringspaces=false,
  tabsize=4,
  morekeywords={include, define},
  % lettere accentate nei commenti
  literate={à}{{\`a}}1 {è}{{\`e}}1 {é}{{\'e}}1 {ì}{{\`i}}1
           {ò}{{\`o}}1 {ù}{{\`u}}1 {È}{{\`E}}1 {À}{{\`A}}1
}
\lstset{style=c}
\newtcblisting{codice}{listing only, nota, left=6mm,
  listing options={style=c}}
\newtcblisting{comando}{listing only, nota,
  listing options={basicstyle=\ttfamily\small, numbers=none, language={},
    columns=flexible, keepspaces=true, breaklines=true}}
\newtcblisting{risultato}{listing only, enhanced jigsaw, breakable,
  colback=black!5, colframe=black, boxrule=0.5pt, sharp corners,
  listing options={basicstyle=\ttfamily\small, numbers=none, language={},
    columns=flexible, keepspaces=true, breaklines=true}}

% --- Diagrammi di flusso (simboli standard) ---
\tikzset{
  terminale/.style = {ellipse, draw, minimum width=2.5cm, minimum height=0.9cm},
  io/.style        = {trapezium, trapezium left angle=70, trapezium right angle=110,
                      draw, minimum width=2.5cm, minimum height=0.9cm},
  processo/.style  = {rectangle, draw, minimum width=2.5cm, minimum height=0.9cm},
  decisione/.style = {diamond, draw, aspect=2, minimum width=2.5cm},
  freccia/.style   = {thick, -{Stealth}}
}

\title{Intorni e limiti all'infinito}
\author{Marco Cerbella}
\date{Lezione 9 -- 22 ottobre 2026}

\begin{document}
\maketitle

\section{Intorni}

\begin{definizione}[Intorno]
Dato un punto $x_0 \in \mathbb{R}$, un \emph{intorno di ampiezza $\varepsilon > 0$} di $x_0$ è un intervallo aperto di lunghezza $2\varepsilon$ che contiene $x_0$, cioè
\[
I_\varepsilon(x_0) = \,]x_0 - \varepsilon,\ x_0 + \varepsilon[ .
\]
\end{definizione}

La \emph{famiglia degli intorni di $x_0$} è l'insieme degli intorni del punto $x_0$ ottenuta al variare di $\varepsilon$ in $]0, +\infty)$.

\begin{proprieta}[degli intorni]
\begin{itemize}
    \item Se $r_1 < r_2$ allora $I_{r_1}(x_0) \subset I_{r_2}(x_0)$.
    \item Dati $x_1, x_2 \in \mathbb{R}$, $x_1 \neq x_2$, esiste sempre $r > 0$ tale che $I_r(x_1) \cap I_r(x_2) = \emptyset$ (basta che $r \le \tfrac{1}{2}|x_1 - x_2|$).
\end{itemize}
\end{proprieta}

\section{Limiti di funzioni per $x \to \pm\infty$}

\begin{definizione}[Limite finito per $x \to +\infty$]
Si dice che $f : D \to \mathbb{R}$ \emph{tende al limite finito} $\ell \in \mathbb{R}$ per $x$ tendente a $+\infty$ e si scrive $\displaystyle\lim_{x \to +\infty} f(x) = \ell$ se
\[
\forall \varepsilon > 0,\ \exists K > 0 \text{ tale che } \forall x \in D : x > K \text{ si ha } |f(x) - \ell| < \varepsilon .
\]
\end{definizione}

La condizione $|f(x) - \ell| < \varepsilon$ significa $f(x) \in I_\varepsilon(\ell)$, cioè $\ell - \varepsilon < f(x) < \ell + \varepsilon$: da un certo punto in poi il grafico di $f$ sta nella striscia di ampiezza $2\varepsilon$ attorno alla retta $y = \ell$.

\begin{definizione}[Limite $+\infty$ per $x \to +\infty$]
Si dice che $f : D \to \mathbb{R}$ \emph{tende a $+\infty$} per $x$ tendente a $+\infty$ e si scrive $\displaystyle\lim_{x \to +\infty} f(x) = +\infty$ se
\[
\forall M > 0,\ \exists K > 0 \text{ tale che } \forall x \in D : x > K \text{ si ha } f(x) > M .
\]
\end{definizione}

\begin{definizione}[Limiti per $x \to -\infty$]
Si dice che $f : D \to \mathbb{R}$ \emph{tende al limite finito} $\ell \in \mathbb{R}$ per $x$ tendente a $-\infty$ e si scrive $\displaystyle\lim_{x \to -\infty} f(x) = \ell$ se
\[
\forall \varepsilon > 0,\ \exists K > 0 \text{ tale che } \forall x \in D : x < -K \text{ si ha } |f(x) - \ell| < \varepsilon .
\]
Si dice che $f$ \emph{tende a $-\infty$} per $x$ tendente a $-\infty$ e si scrive $\displaystyle\lim_{x \to -\infty} f(x) = -\infty$ se
\[
\forall M > 0,\ \exists K > 0 \text{ tale che } \forall x \in D : x < -K \text{ si ha } f(x) < -M .
\]
\end{definizione}

\begin{osservazione}
Nelle slide la condizione per $x \to -\infty$ è scritta $x < K$ e $f(x) < M$; nelle verifiche svolte a lezione si usa però, come sopra, $x < -K$ (e $f(x) < -M$ per il limite $-\infty$), con $K, M > 0$.
\end{osservazione}

Il caso $\lim_{x\to-\infty} f(x) = +\infty$ e quello $\lim_{x\to+\infty} f(x) = -\infty$ si definiscono in modo analogo, combinando le condizioni su $x$ e su $f(x)$.

\section{Verifiche con la definizione}

\begin{esempio}[$\lim_{x\to+\infty}\frac{x^2+2x}{2x^2+1} = \frac{1}{2}$]
Il campo di esistenza: $2x^2 + 1 \neq 0$ per ogni $x \in \mathbb{R}$, quindi $D = \mathbb{R}$ e $x$ può tendere a $+\infty$ perché $D$ non è limitato superiormente.

Fissato $\varepsilon > 0$, la condizione $|f(x) - \ell| < \varepsilon$ si traduce in
\[
\left| \frac{x^2 + 2x}{2x^2 + 1} - \frac{1}{2} \right| < \varepsilon .
\]
Calcoliamo la differenza:
\[
\frac{x^2 + 2x}{2x^2 + 1} - \frac{1}{2} = \frac{2x^2 + 4x - 2x^2 - 1}{2(2x^2 + 1)} = \frac{4x - 1}{2(2x^2 + 1)} .
\]
Supponiamo $x > \tfrac{1}{4}$, così togliamo il valore assoluto:
\[
\frac{4x - 1}{2(2x^2 + 1)} < \frac{4x}{2(2x^2 + 1)} = \frac{2x}{2x^2 + 1} < \frac{2x}{2x^2} = \frac{1}{x} .
\]
Se $\dfrac{1}{x} < \varepsilon$, cioè $x > \dfrac{1}{\varepsilon}$, la disuguaglianza è verificata. Basta quindi prendere
\[
K = \max\left\{ \tfrac{1}{4},\ \tfrac{1}{\varepsilon} \right\},
\]
perché se $x > K$ allora $x > \tfrac{1}{4}$ e $x > \tfrac{1}{\varepsilon}$, e quindi $\left|\dfrac{4x-1}{2(2x^2+1)}\right| < \varepsilon$.
\end{esempio}

\begin{esempio}[$\lim_{x\to+\infty}\sqrt{x} = +\infty$]
Il campo di esistenza è $x \ge 0$, cioè $D = [0, +\infty)$. Fissato $M > 0$, la condizione $f(x) > M$ si traduce in $\sqrt{x} > M$, cioè $x > M^2$. Basta considerare $K = M^2$: per ogni $M > 0$, considerando $K = M^2$, per ogni $x > M^2$ si ha $\sqrt{x} > M$.
\end{esempio}

\begin{esempio}[$\lim_{x\to-\infty}\frac{1}{\sqrt{1-x}} = 0$]
Campo di esistenza: $\sqrt{1-x} \neq 0$ e $1 - x \ge 0$, cioè $x \ne 1$ e $x \le 1$, quindi $x < 1$ e $D = (-\infty, 1)$.

Fissato $\varepsilon > 0$,
\[
|f(x) - \ell| = \left| \frac{1}{\sqrt{1-x}} \right| = \frac{1}{\sqrt{1-x}} < \varepsilon
\iff \sqrt{1-x} > \frac{1}{\varepsilon}
\iff 1 - x > \frac{1}{\varepsilon^2}
\iff x < 1 - \frac{1}{\varepsilon^2} .
\]
Si prende $K = \max\left\{0,\ \dfrac{1}{\varepsilon^2} - 1\right\}$: per ogni $x < -K$ si ha $\dfrac{1}{\sqrt{1-x}} < \varepsilon$.
\end{esempio}

\end{document}
